\chapter{Compensation et règlement-livraison}\label{ch:m1:clearing-and-settlement}

Le 28 janvier 2021, à 5 h 11 du matin, un avis automatique parvint
à la filiale de \omterm{def:m1:clearing-and-settlement:clearing}{compensation} d'un courtier de détail. Pendant la nuit, le collatéral que
la chambre de \omterm{def:m1:clearing-and-settlement:clearing}{compensation} des actions américaines lui réclamait avait grimpé à environ 3.7~milliards de dollars. La société
disposait d'environ 700~millions de dollars en dépôt~; les 3~milliards de dollars restants devaient
être réglés avant dix heures. Aucune transaction n'avait posé problème. Ses clients avaient acheté
d'énormes quantités de quelques actions aux cours très volatils, ces achats ne
seraient réglés que deux jours plus tard, et d'ici là quelqu'un devait
se porter garant de ces achats. Ce tiers fait l'objet de ce chapitre~: ce qui se passe
entre le moment où une transaction est convenue et celui où elle devient définitive, qui
la garantit dans l'intervalle, et ce que coûte cette garantie.

\section{De la transaction au règlement-livraison}

\begin{definition}[Compensation et règlement-livraison]\label{def:m1:clearing-and-settlement:clearing}
Le \emph{règlement-livraison}\index{règlement-livraison} est l'échange définitif de ce qui a été négocié~:
les titres passent à l'acheteur, les espèces au vendeur, et la transaction ne peut plus
être annulée. La \emph{compensation}\index{compensation} désigne tout ce qui se situe entre l'exécution
et le règlement-livraison~: confirmation des termes, calcul de ce que chacun doit à qui,
le netting, et gestion du risque qu'une partie fasse défaut avant d'avoir payé ou
livré.
\end{definition}

\begin{definition}[Cycle de règlement-livraison et suspens de règlement-livraison]\label{def:m1:clearing-and-settlement:cycle}
Le \emph{cycle de règlement-livraison}\index{cycle de règlement-livraison} est le délai standard
entre la date de transaction $T$ et la date de \omterm{def:m1:clearing-and-settlement:clearing}{règlement-livraison}, noté $T{+}n$ pour $n$ jours
ouvrés. Un \emph{suspens de règlement-livraison}\index{suspens de règlement-livraison} est une obligation non
honorée à sa date de \omterm{def:m1:clearing-and-settlement:clearing}{règlement-livraison}, généralement parce que le vendeur ne disposait pas des
titres à livrer.
\end{definition}

\begin{dated}{2026-09}{Cycles de règlement-livraison}\label{dat:m1:clearing-and-settlement:tplus1}
Les titres américains se règlent à $T{+}1$ depuis le 28 mai 2024. L'Union européenne
a légiféré dans le même sens pour le 11 octobre 2027 (l'amendement à son
règlement sur les dépositaires centraux de titres, publié en octobre 2025)~;
le Royaume-Uni et la Suisse se sont alignés sur la même date. Jusque-là
la plupart des titres européens se règlent à $T{+}2$.
\end{dated}

\begin{definition}[Dépositaire central de titres et livraison contre paiement]\label{def:m1:clearing-and-settlement:csd}
Un \emph{dépositaire central de titres}\index{dépositaire central de titres}
(CSD) tient le registre définitif des propriétaires de chaque titre d'un marché et
assure le règlement des transactions par jeu d'écritures entre les comptes de ses participants.
La \emph{livraison contre paiement}\index{livraison contre paiement} (DVP) est la
règle selon laquelle la jambe titres et la jambe espèces d'un règlement deviennent
définitives ensemble ou pas du tout, de sorte qu'aucune des deux parties ne puisse perdre
la valeur totale de la transaction.
\end{definition}

\begin{omfigure}
\begin{tikzpicture}[font=\small,
  ev/.style={draw=omDef,rounded corners=2pt,minimum width=2.7cm,minimum height=0.85cm,align=center,fill=omDef!4}]
\draw[-{Stealth[length=2.2mm]},thick,omExo] (-0.6,0) -- (15.2,0);
\foreach \x/\lab in {0.6/{$T$, 10:31},4.4/{$T$, soir},8.2/{$T{+}1$, matin},12/{$T{+}1$, journée}}{
  \draw[omExo,thick] (\x,-0.12) -- (\x,0.12);
  \node[font=\footnotesize,omExo,anchor=north] at (\x,-0.2) {\lab};
}
\node[ev] at (0.6,1.05) {Exécution\\sur une plateforme};
\node[ev] at (4.4,1.05) {Novation,\\compensation};
\node[ev] at (8.2,1.05) {Appel de marge\\honoré};
\node[ev] at (12,1.05) {Règlement-livraison\\DVP au CSD};
\draw[omThm,thick,decorate,decoration={brace,mirror,amplitude=5pt}] (0.6,-0.95) -- (12,-0.95)
  node[midway,below=6pt,font=\footnotesize,omThm] {la chambre de compensation est exposée au défaut de chaque membre};
\end{tikzpicture}

\omcaption{Le parcours d'une transaction sur action sous un cycle d'un jour. L'accolade est la
fenêtre dont traite la suite du chapitre~: sa durée, multipliée par la volatilité de
ce qui a été négocié, fixe le collatéral.}\label{fig:m1:clearing-and-settlement:timeline}
\end{omfigure}

\section{La contrepartie centrale}

\begin{definition}[Contrepartie centrale et novation]\label{def:m1:clearing-and-settlement:ccp}
Une \emph{contrepartie centrale}\index{contrepartie centrale} (CCP), ou chambre de
\omterm{def:m1:clearing-and-settlement:clearing}{compensation}, s'interpose dans chaque transaction qu'elle accepte~: par
\emph{novation}\index{novation}, le contrat entre acheteur et vendeur est
remplacé par deux contrats, l'un entre l'acheteur et la CCP et l'autre entre
la CCP et le vendeur. Chaque membre n'a alors qu'une seule contrepartie, quel que soit
le nombre de sociétés avec lesquelles il a négocié.
\end{definition}

\begin{proposition}[Effet de la novation sur un réseau]\label{prop:m1:clearing-and-settlement:network}
Entre $n$ membres négociant de gré à gré, il peut y avoir jusqu'à $n(n-1)/2$ relations
de contrepartie~; avec une CCP, il n'en existe que $n$. Après la \omterm{def:m1:clearing-and-settlement:ccp}{novation}, les obligations de chaque membre
pour un même titre et une même devise se réduisent à une position nette unique face à
la CCP, et ces positions nettes ont une somme nulle.
\end{proposition}
\begin{proof}
Une paire de membres est un sous-ensemble de deux éléments parmi $n$~: $\binom n2$ paires. Avec une CCP
chaque contrat a la CCP d'un côté, ce qui donne une relation par membre.
Les contrats d'un membre avec la CCP sur un même titre sont fongibles et se cumulent en
$\sum(\text{achats}) - \sum(\text{ventes})$. Chaque transaction ajoute $+q$ à un
membre et $-q$ à un autre, de sorte que la somme sur les membres est nulle.
\end{proof}

\begin{omfigure}
\begin{tikzpicture}[font=\small,
  m/.style={circle,draw=omDef,fill=omDef!6,minimum size=0.75cm,inner sep=0pt},
  c/.style={circle,draw=omThm,fill=omThm!8,minimum size=1.05cm,inner sep=0pt}]
\begin{scope}
\foreach \i/\a in {A/90,B/30,C/330,D/270,E/210,F/150}{ \node[m] (\i) at (\a:1.9) {\i}; }
\foreach \i/\j in {A/B,A/C,A/D,A/E,A/F,B/C,B/D,B/E,B/F,C/D,C/E,C/F,D/E,D/F,E/F}{ \draw[omExo] (\i) -- (\j); }
\node[font=\footnotesize] at (0,-2.75) {bilatéral~: 15 relations};
\end{scope}
\begin{scope}[xshift=7.6cm]
\node[c] (K) at (0,0) {CCP};
\foreach \i/\a in {A/90,B/30,C/330,D/270,E/210,F/150}{ \node[m] (\i) at (\a:1.9) {\i}; \draw[omDef,thick] (\i) -- (K); }
\node[font=\footnotesize] at (0,-2.75) {compensé~: 6 relations};
\end{scope}
\end{tikzpicture}

\omcaption{Six membres avant et après la \omterm{def:m1:clearing-and-settlement:ccp}{novation}. La \omterm{def:m1:clearing-and-settlement:ccp}{contrepartie centrale} supprime le réseau,
et concentre en un seul nœud tout le risque qui était réparti sur lui.}
\end{omfigure}

\begin{example}[Le netting des transactions d'un matin]\label{ex:m1:clearing-and-settlement:netting}
A achète 5\,000 actions à B au prix de 20.00~\$, B achète 3\,000 à C au prix de 20.10~\$, et C
achète 4\,000 à A au prix de 19.95~\$. Brut, 12\,000 actions et 240\,100~\$
circuleraient. En net, A reçoit 1\,000 actions et paie 20\,200~\$~; B livre 2\,000 et
reçoit 39\,700~\$~; C reçoit 1\,000 et paie 19\,500~\$. Trois livraisons
totalisant 2\,000 actions et 39\,700~\$ sont réglées à la place~: le netting a supprimé
83~\% de la valeur. Le tutoriel montre que le ratio monte vers 98~\% à mesure que
l'activité augmente, ce qui explique qu'un marché de milliers de membres puisse se
régler tout simplement.
\end{example}

\section{Marges et cascade de défaillance}

Une \omterm{def:m1:clearing-and-settlement:ccp}{contrepartie centrale} qui garantit chaque transaction doit survivre à la défaillance
d'un membre. Elle y parvient avec l'argent d'autrui, appelé dans un ordre fixe.

\begin{definition}[Marge de variation et marge initiale]\label{def:m1:clearing-and-settlement:margin}
La \emph{marge de variation}\index{marge de variation} est le paiement quotidien (ou intrajournalier)
du changement de valeur des positions ouvertes d'un membre~: les perdants paient, la
\omterm{def:m1:clearing-and-settlement:ccp}{contrepartie centrale} reverse les liquidités aux gagnants, et les expositions repartent de zéro.
La \emph{marge initiale}\index{marge initiale} est un collatéral déposé pour couvrir
la perte que la \omterm{def:m1:clearing-and-settlement:ccp}{contrepartie centrale} pourrait subir sur les positions du membre entre son
dernier paiement de marge de variation et le moment où la \omterm{def:m1:clearing-and-settlement:ccp}{contrepartie centrale} a fini de
les clôturer, un délai appelé période de marge en risque.
\end{definition}

\begin{method}[Une marge initiale en valeur à risque]\label{met:m1:clearing-and-settlement:im}
Pour une position nette de valeur $N$ sur un instrument de volatilité quotidienne
$\sigma$, à couvrir sur $d$ jours avec un niveau de confiance $\Phi(z)$~:
\[
\mathrm{IM} = z\,\sigma\sqrt{d}\;|N| .
\]
Les \omterm{def:m1:clearing-and-settlement:ccp}{contreparties centrales} ajoutent des majorations pour la concentration, pour l'illiquidité, pour les écarts entre
régimes de volatilité et --- dans la chambre de \omterm{def:m1:clearing-and-settlement:clearing}{compensation} américaine des actions --- pour un membre
dont l'exigence est élevée par rapport à ses propres fonds propres. Les modèles réels sont
fondés sur le portefeuille~; le \cref{ch:m1:margin} en construit un.
\end{method}

\begin{example}[Pourquoi l'appel de marge est tombé]\label{ex:m1:clearing-and-settlement:why}
Appliquer la méthode à un courtier dont les clients ont acheté, en net, pour
1~milliard de dollars d'une action, avec $\sigma = 4\%$, $d=2$, $z = 2.33$, $\mathrm{IM} = 132$ millions de dollars.
Si la volatilité de l'action passe à 25~\% par jour, la même position exige
824~millions de dollars, et doubler la position la double à nouveau. En janvier 2021
la composante liée à la volatilité de l'exigence du courtier de détail était d'environ
1.3~milliard de dollars, à laquelle la formule de la chambre de \omterm{def:m1:clearing-and-settlement:clearing}{compensation} ajoutait une «~prime
de capital excédentaire~» de 2.2~milliards de dollars, parce que l'exigence dépassait largement les fonds propres de la société.
La chambre de \omterm{def:m1:clearing-and-settlement:clearing}{compensation} a renoncé à cette prime le matin même, pour cette
société et d'autres~: 9.7~milliards de dollars au total pour l'ensemble de ses membres cette semaine-là,
selon l'enquête du Congrès américain. Le courtier a restreint
les achats sur les actions concernées --- seul moyen dont il disposait pour empêcher la
hausse de l'exigence --- et entreprit, la même semaine, de lever
3.5~milliards de dollars de nouveaux fonds propres. Raccourcir le \omterm{def:m1:clearing-and-settlement:cycle}{cycle de règlement-livraison} à $T{+}1$ divise une telle
exigence par $\sqrt2$~: c'était là, et non la commodité, l'argument en sa faveur.
\end{example}

\begin{omfigure}
\begin{tikzpicture}
\begin{axis}[width=10.5cm,height=4.8cm,
  xlabel={volatilité quotidienne de la position (\%)},ylabel={marge (\% de la position)},
  tick label style={font=\small},label style={font=\small},
  xmin=1,xmax=15,ymin=0,ymax=55,grid=major,grid style={gray!25},
  legend columns=2,legend style={font=\footnotesize,at={(0.5,-0.34)},anchor=north,draw=none,fill=none,/tikz/every even column/.append style={column sep=8pt}},
  table/col sep=comma]
\addplot[omThm,very thick] table[x=sigma_pct,y=t2]{figdata/markets-1/05-clearing-and-settlement/margin_vs_vol.csv};
\addplot[omDef,very thick] table[x=sigma_pct,y=t1]{figdata/markets-1/05-clearing-and-settlement/margin_vs_vol.csv};
\legend{exposition à deux jours,exposition à un jour}
\end{axis}
\end{tikzpicture}

\omcaption{La marge en valeur à risque de la \cref{met:m1:clearing-and-settlement:im}
à 99~\% de confiance. La marge est linéaire en volatilité~: une action qui devient
cinq fois plus volatile du jour au lendemain coûte cinq fois plus cher à compenser le
matin suivant. Données~: calculées par le script du chapitre.}
\end{omfigure}

\begin{definition}[Fonds de défaillance et cascade de défaillance]\label{def:m1:clearing-and-settlement:waterfall}
Le \emph{fonds de défaillance}\index{fonds de défaillance} est une réserve de collatéral
apportée par tous les membres pour absorber les pertes qui excèdent les ressources propres
du défaillant. La \emph{cascade de défaillance}\index{cascade de défaillance} est l'ordre
dans lequel les ressources absorbent la perte causée par la défaillance d'un membre~:
\begin{enumerate}
\item la \omterm{def:m1:clearing-and-settlement:margin}{marge initiale} du défaillant~;
\item la contribution du défaillant au fonds de défaillance~;
\item une tranche des fonds propres de la \omterm{def:m1:clearing-and-settlement:ccp}{contrepartie centrale} (sa «~part de risque~»)~;
\item les contributions des membres survivants au fonds de défaillance~;
\item des appels complémentaires que la \omterm{def:m1:clearing-and-settlement:ccp}{contrepartie centrale} peut réclamer aux survivants, jusqu'à un plafond.
\end{enumerate}
\end{definition}

La conception tient à ce principe «~le défaillant paie d'abord, les survivants ensuite, tous ensemble~»~: les
couches 1 et 2 font répondre chaque membre de son propre risque~; la couche 3 donne aux
actionnaires de la \omterm{def:m1:clearing-and-settlement:ccp}{contrepartie centrale} une raison de fixer correctement les marges~; les couches 4 et 5 font de chaque membre un
garant de tous les autres --- ce qui explique que les membres se soucient de qui est admis
et de ce qu'il leur est permis de faire.

\begin{omfigure}
\begin{tikzpicture}
\begin{axis}[width=11cm,height=5.4cm,stack plots=y,area style,no markers,
  xlabel={perte liée à la défaillance (millions de \$)},ylabel={qui paie (millions de \$)},
  tick label style={font=\small},label style={font=\small},
  xmin=0,xmax=1100,ymin=0,ymax=1100,
  legend columns=3,legend cell align=left,legend style={font=\footnotesize,at={(0.5,-0.3)},anchor=north,draw=none,fill=none,/tikz/every even column/.append style={column sep=6pt}},
  table/col sep=comma]
\addplot[fill=omDef!70,draw=omDef] table[x=loss,y=margin]{figdata/markets-1/05-clearing-and-settlement/waterfall.csv} \closedcycle;
\addplot[fill=omDef!35,draw=omDef] table[x=loss,y=own_fund]{figdata/markets-1/05-clearing-and-settlement/waterfall.csv} \closedcycle;
\addplot[fill=omMeth!55,draw=omMeth] table[x=loss,y=ccp]{figdata/markets-1/05-clearing-and-settlement/waterfall.csv} \closedcycle;
\addplot[fill=omProp!55,draw=omProp] table[x=loss,y=survivors]{figdata/markets-1/05-clearing-and-settlement/waterfall.csv} \closedcycle;
\addplot[fill=omThm!45,draw=omThm] table[x=loss,y=assess]{figdata/markets-1/05-clearing-and-settlement/waterfall.csv} \closedcycle;
\addplot[fill=black!55,draw=black] table[x=loss,y=uncovered]{figdata/markets-1/05-clearing-and-settlement/waterfall.csv} \closedcycle;
\legend{marge du défaillant,part du défaillant au fonds,capital de la CCP,fonds des survivants,appels complémentaires,non couvert}
\end{axis}
\end{tikzpicture}

\omcaption{Qui absorbe une perte de défaut à mesure qu'elle croît, pour des couches illustratives de
120, 30, 20, 400 et 400 millions. Jusqu'à 150, le défaillant paie lui-même~;
à partir de 170, ce sont les autres membres qui paient~; au-delà de 970, personne n'est engagé. Données~:
calculées par le script du chapitre.}\label{fig:m1:clearing-and-settlement:waterfall}
\end{omfigure}

\section{Lorsqu'un membre fait défaut}

\begin{method}[La séquence de gestion du défaut]\label{met:m1:clearing-and-settlement:dm}
\begin{enumerate}
\item \textbf{Déclarer} le défaut et cesser d'accepter les transactions du membre.
\item \textbf{Transférer} les clients du défaillant, avec leurs positions et
  leur marge, vers d'autres membres lorsque la structure des comptes le permet.
\item \textbf{Couvrir} le portefeuille propre du défaillant afin de neutraliser son
  risque de marché dans les heures qui suivent.
\item \textbf{Mettre aux enchères} le portefeuille couvert auprès des membres survivants, qui
  sont tenus ou fortement incités à enchérir.
\item \textbf{Répartir} toute perte le long de la \omterm{def:m1:clearing-and-settlement:waterfall}{cascade de défaillance}~; restituer ce qui
  reste de la marge du défaillant à son administrateur.
\end{enumerate}
\end{method}

\begin{example}[Septembre 2008]\label{ex:m1:clearing-and-settlement:lehman}
Lorsque Lehman Brothers fait défaut le lundi 15 septembre 2008, la chambre de \omterm{def:m1:clearing-and-settlement:clearing}{compensation} londonienne
des swaps de taux d'intérêt détient son portefeuille~: 66\,390 transactions
d'un notionnel de 9~billions de dollars répartis sur cinq devises, face à environ
2~milliards de dollars de \omterm{def:m1:clearing-and-settlement:margin}{marge initiale}. Des traders détachés par les banques membres couvrent le
portefeuille aux côtés de l'équipe des risques de la chambre de \omterm{def:m1:clearing-and-settlement:clearing}{compensation}~; entre le 24 septembre et
le 3 octobre, les portefeuilles couverts par devise sont mis aux enchères. La chambre de \omterm{def:m1:clearing-and-settlement:clearing}{compensation}
indique que le défaut a été géré largement dans les limites de la marge détenue et que
son \omterm{def:m1:clearing-and-settlement:waterfall}{fonds de défaillance} n'a pas été utilisé. L'épisode est devenu l'argument de référence
en faveur des obligations de \omterm{def:m1:clearing-and-settlement:clearing}{compensation} qui ont suivi.
\end{example}

\begin{remark}[Le risque n'a pas disparu]\label{rem:m1:clearing-and-settlement:concentration}
Une \omterm{def:m1:clearing-and-settlement:ccp}{contrepartie centrale} n'élimine pas le risque de contrepartie~; elle remplace de nombreuses
petites expositions opaques par une seule, grande, collatéralisée et supervisée. Ses appels de
marge sont \emph{procycliques}~: ils augmentent précisément lorsque la volatilité est élevée et que
la liquidité se raréfie, de sorte que le mécanisme qui protège le système peut assécher la
liquidité de ses membres au pire moment --- le thème des
\cref{ch:m1:margin,ch:m1:stress-case-studies}. Et la chambre de \omterm{def:m1:clearing-and-settlement:clearing}{compensation} elle-même
devient une entreprise qui ne peut pas être autorisée à faire défaut.
\end{remark}

\section{Tutoriel~: compenser une journée de transactions}\label{tut:m1:clearing-and-settlement:netting}

\begin{tutorial}
\textbf{Objectif.} Compenser une journée de transactions de façon multilatérale et mesurer ce que la \omterm{def:m1:clearing-and-settlement:clearing}{compensation}
supprime. \textbf{État final~:} le graphique ci-dessous, de 41~\% de valeur supprimée pour
dix transactions à 98~\% pour dix mille.
\begin{tutsteps}
\item \textbf{Compenser.} Un parcours unique des transactions accumule, par membre, les actions
par symbole et le cash.
\omcode{code/markets-1/05-clearing-and-settlement/python/clearing_demo.py}{18}{35}{Compensation multilatérale et part de valeur supprimée.}\label{lst:m1:clearing-and-settlement:net}
\item \textbf{Vérifier les invariants.} Le cash net s'annule sur l'ensemble des membres,
de même pour les actions nettes dans chaque symbole~: les tests vérifient les deux. Un moteur
de \omterm{def:m1:clearing-and-settlement:clearing}{compensation} qui les viole a perdu une transaction.
\item \textbf{Mesurer.} Générer des transactions aléatoires entre huit membres sur
cinq actions et moyenner l'efficacité sur vingt graines.
\item \textbf{Allouer une perte.} La cascade tient en six lignes.
\omcode{code/markets-1/05-clearing-and-settlement/python/clearing_demo.py}{58}{65}{Faire descendre une perte le long de la cascade.}\label{lst:m1:clearing-and-settlement:wf}
\end{tutsteps}
\textbf{Pour aller plus loin.} Faire d'un membre un courtier de détail qui n'achète qu'une
seule action~: son efficacité de \omterm{def:m1:clearing-and-settlement:clearing}{compensation} s'effondre, et avec elle l'économie de marge que
la \omterm{def:m1:clearing-and-settlement:clearing}{compensation} procure aux autres. Attribuer ensuite à chaque membre une \omterm{def:m1:clearing-and-settlement:margin}{marge initiale},
selon la \cref{met:m1:clearing-and-settlement:im}, sur sa position nette, et repérer le membre
dont l'exigence est la plus élevée par rapport à sa taille.
\end{tutorial}

\begin{omfigure}
\begin{tikzpicture}
\begin{axis}[width=10.5cm,height=4.6cm,xmode=log,
  xlabel={transactions dans la journée (huit membres, cinq actions)},ylabel={valeur compensée (\%)},
  tick label style={font=\small},label style={font=\small},
  xmin=10,xmax=10000,ymin=30,ymax=100,grid=major,grid style={gray!25},table/col sep=comma]
\addplot[omDef,very thick,mark=*,mark size=1.4pt] table[x=trades,y=eff_pct]{figdata/markets-1/05-clearing-and-settlement/netting.csv};
\end{axis}
\end{tikzpicture}

\omcaption{Efficacité de la \omterm{def:m1:clearing-and-settlement:clearing}{compensation} multilatérale selon l'activité (axe horizontal
logarithmique). Plus un marché connaît de flux dans les deux sens, plus la
fraction qui doit jamais être réglée est petite. Données~: la simulation du
tutoriel, moyenne de vingt graines par point.}
\end{omfigure}

\section{Construction~: le moteur de compensation}\label{bld:m1:clearing-and-settlement:netting}

\begin{build}
\bfield{Rôle} La mini-société compense son propre marché simulé. À la
fin de chaque session, elle doit savoir, par membre, ce qui est réglé.
\bfield{Interface} \texttt{NettingEngine.add(trade)} pendant la session~;
\texttt{close()} renvoyant une structure immuable \texttt{SettlementInstructions}~:
par membre et par symbole une quantité d'actions signée, par membre et par devise un
montant de cash signé en unités de registre (\cref{ch:m1:what-a-trading-firm-does}),
et la liste des identifiants de transaction compensés dans chacun.
\bfield{Règles} Les prix et les montants de cash sont des entiers en unités de registre~: pas de virgule
flottante. Une transaction dont l'acheteur est égal au vendeur, une quantité non positive ou un
identifiant en double est rejetée. Après \texttt{close()}, aucune transaction n'est
acceptée. Les deux invariants à somme nulle sont vérifiés à l'intérieur de \texttt{close()}
et leur violation est une erreur fatale, non un avertissement.
\bfield{Tests d'acceptation} \texttt{code/firm/clearing/tests/}~: reproduit
l'\cref{ex:m1:clearing-and-settlement:netting} à l'identique~; invariants sur dix
mille transactions aléatoires~; rejets~; chaque identifiant de transaction figure dans exactement
deux jeux d'instructions de membres.
\bfield{Pour aller plus loin} Ajouter \texttt{fail(member, symbol, quantity)}~: une livraison partielle,
le suspens qui en résulte reporté au lendemain, et le buy-in qui
le clôt.
\end{build}

\begin{omsources}
\item US House Committee on Financial Services, rapport du personnel majoritaire,
\emph{Game Stopped: How the Meme Stock Market Event Exposed Troubling
Business Practices, Inadequate Risk Management, and the Need for Legislative
and Regulatory Reform}, juin 2022.
\item US Securities and Exchange Commission, \emph{Staff Report on Equity and
Options Market Structure Conditions in Early 2021}, octobre 2021.
\item US Securities and Exchange Commission, communiqué de presse 2024-62 sur le
$T{+}1$ cycle de règlement-livraison; Euronext, \emph{T+1 programme}; UK Financial
Conduct Authority, \emph{About T+1 settlement}.
\item LCH.Clearnet, \emph{\$9 trillion Lehman OTC interest rate swap default
successfully resolved}, communiqué de presse, 8 octobre 2008.
\item Committee on Payments and Market Infrastructures and IOSCO,
\emph{Principles for financial market infrastructures}, avril 2012.
\item J.~Gregory, \emph{Central Counterparties}, Wiley, 2014.
\end{omsources}

\section{Exercices}

\begin{exercise}[$\star$]\label{exo:m1:clearing-and-settlement:1}
Combien de relations bilatérales peuvent exister entre 40 membres~? Combien
avec une \omterm{def:m1:clearing-and-settlement:ccp}{contrepartie centrale}~? De quel facteur le nombre diminue-t-il~?
\end{exercise}

\begin{exercise}[$\star$]\label{exo:m1:clearing-and-settlement:2}
A vend 2\,000 actions à B à 50.00~\$~; B vend 2\,000 actions à C à 50.20~\$~; C vend
1\,500 actions à A à 50.10~\$. Donner les actions nettes et les flux de trésorerie nets de chaque membre, ainsi que
l'efficacité du netting en valeur.
\end{exercise}

\begin{exercise}[$\star$]\label{exo:m1:clearing-and-settlement:3}
Calculer la \omterm{def:m1:clearing-and-settlement:margin}{marge initiale} fournie par la \cref{met:m1:clearing-and-settlement:im} sur
une position nette de 400~millions de dollars, avec $\sigma = 2.5\%$, $z = 2.33$, pour une
période de deux jours et pour une période d'un jour. Que libère le cycle le plus court~?
\end{exercise}

\begin{exercise}[$\star\star$]\label{exo:m1:clearing-and-settlement:4}
Avec la \omterm{def:m1:clearing-and-settlement:waterfall}{cascade de défaillance} de la \cref{fig:m1:clearing-and-settlement:waterfall},
répartir les pertes de (a) 140~millions de dollars, (b) 300~millions de dollars, (c) 1\,000
millions de dollars. Il reste 25 membres survivants, dont les contributions au \omterm{def:m1:clearing-and-settlement:waterfall}{fonds de défaillance} sont égales~: que
perd chacun en (b)~?
\end{exercise}

\begin{exercise}[$\star\star$]\label{exo:m1:clearing-and-settlement:5}
Les clients d'un courtier détiennent une position longue nette non réglée de 600~millions de dollars.
Pendant la nuit, la volatilité quotidienne du titre est réestimée, passant de 3~\% à 18~\%.
Calculer la marge avant et après ce changement ($z = 2.33$, deux jours). Le
capital du courtier est de 500~millions de dollars~: commenter.
\end{exercise}

\begin{exercise}[$\star\star$]\label{exo:m1:clearing-and-settlement:6}
Sans \omterm{def:m1:clearing-and-settlement:csd}{livraison contre paiement} (DVP), un vendeur livre 10~millions de dollars d'actions le matin et doit
être payé l'après-midi. Avec la \omterm{def:m1:clearing-and-settlement:csd}{livraison contre paiement} et une \omterm{def:m1:clearing-and-settlement:ccp}{contrepartie centrale}, la transaction est garantie à
son prix. Si l'acheteur fait défaut à midi, après une baisse de 4~\% du titre, que
perd le vendeur dans chacun des cas~?
\end{exercise}

\begin{exercise}[$\star\star\star$]\label{exo:m1:clearing-and-settlement:7}
\emph{Programmation.} En utilisant \texttt{random\_trades(2000, 8, seed=1)}, calculer
l'efficacité du netting. Ajouter ensuite 400 transactions dans lesquelles le membre \texttt{A} achète
500 actions de \texttt{S0} auprès d'un autre membre tiré au hasard, puis recalculer
l'efficacité pour l'ensemble des transactions. Expliquer le signe de la variation.
\end{exercise}

\begin{exercise}[$\star\star\star$]\label{exo:m1:clearing-and-settlement:8}
\emph{Trouver la faille.} «~Notre chambre de \omterm{def:m1:clearing-and-settlement:clearing}{compensation} n'a jamais utilisé son \omterm{def:m1:clearing-and-settlement:waterfall}{fonds de défaillance},
donc ses marges sont manifestement suffisantes~; nous pouvons réduire le fonds et baisser
les coûts des membres~». Donner deux raisons pour lesquelles la prémisse ne soutient pas la
conclusion~: l'une statistique, l'autre liée au comportement.
\end{exercise}

\section{Problème~: le défaut du membre C}

\begin{problem}[{Problème du week-end --- une chambre de compensation face à son pire lundi}]\label{pb:m1:clearing-and-settlement:1}
Une chambre de \omterm{def:m1:clearing-and-settlement:clearing}{compensation} pour les contrats à terme sur indices actions compte 21 membres. Le membre C, une
\omterm{def:m1:what-a-trading-firm-does:prop}{société de trading pour compte propre}, est en position longue sur des contrats à terme pour un notionnel de 2.4~milliards de dollars. La
chambre de \omterm{def:m1:clearing-and-settlement:clearing}{compensation} détient de C une \omterm{def:m1:clearing-and-settlement:margin}{marge initiale} de 130~millions de dollars et une
contribution de 25~millions de dollars au \omterm{def:m1:clearing-and-settlement:waterfall}{fonds de défaillance}. Sa propre tranche de fonds propres est de 15
millions de dollars~; les 20 autres membres apportent chacun 20~millions de dollars au
\omterm{def:m1:clearing-and-settlement:waterfall}{fonds de défaillance} et peuvent chacun être sollicités une fois de plus pour le même montant.

\medskip\noindent\textbf{Partie~I --- La marge était-elle correcte~?}
\begin{enumerate}
\item L'indice a une volatilité quotidienne de 1.3~\%. Quelle marge donne
la \cref{met:m1:clearing-and-settlement:im} pour $d = 2$ jours à 99~\%
($z = 2.33$)~?
\item La chambre de \omterm{def:m1:clearing-and-settlement:clearing}{compensation} détient 130~millions de dollars. À quel niveau de confiance
cela correspond-il~?
\item Le vendredi, l'indice baisse de 3~\%. Quelle \omterm{def:m1:clearing-and-settlement:margin}{marge de variation} est appelée à C~?
\item C la verse. Quelle est l'exposition de la chambre de \omterm{def:m1:clearing-and-settlement:clearing}{compensation} envers C le vendredi
soir~?
\end{enumerate}

\medskip\noindent\textbf{Partie~II --- Le défaut.}
\begin{enumerate}[resume]
\item Le lundi matin, l'indice ouvre en baisse de 6~\% et C ne paie pas.
Quel montant C doit-il~?
\item La chambre de \omterm{def:m1:clearing-and-settlement:clearing}{compensation} déclare C en défaut à 9:00. La couverture du
portefeuille prend jusqu'à midi, et pendant ce temps l'indice baisse encore de 2~\%.
Calculer la perte totale sur la position depuis la clôture du vendredi.
\item La vente aux enchères du portefeuille couvert coûte en outre 0.4~\% du notionnel
en concessions accordées aux enchérisseurs. Donner la perte totale à répartir.
\item Répartir cette perte dans la \omterm{def:m1:clearing-and-settlement:waterfall}{cascade de défaillance}, couche par couche.
\item Quelle perte subit chaque membre survivant~?
\end{enumerate}

\medskip\noindent\textbf{Partie~III --- Scénarios contrefactuels.}
\begin{enumerate}[resume]
\item Si la chambre de \omterm{def:m1:clearing-and-settlement:clearing}{compensation} s'était couverte dans l'heure, avant la baisse
supplémentaire, quelle perte auraient subie les survivants~?
\item Quelle \omterm{def:m1:clearing-and-settlement:margin}{marge initiale} aurait couvert la totalité de la perte de la question~7~?
Exprimer ce montant sous la forme d'un multiple de la volatilité quotidienne multipliée par le notionnel.
\item Quelle est la perte maximale que l'ensemble de la cascade peut absorber~?
\item À quelle baisse totale de l'indice depuis la clôture du vendredi cela correspond-il,
avec le même coût de vente aux enchères~?
\end{enumerate}

\medskip\noindent\textbf{Partie~IV --- Conception.}
\begin{enumerate}[resume]
\item Après l'événement, la chambre de \omterm{def:m1:clearing-and-settlement:clearing}{compensation} porte $z$ à 3 et la période de
marge à trois jours. Donner la nouvelle marge sur une position de 2.4~milliards de dollars.
\item La position de C représentait 30~\% des positions ouvertes. Proposer une
majoration pour concentration et justifier sa forme.
\item Pourquoi les fonds propres de la chambre de \omterm{def:m1:clearing-and-settlement:clearing}{compensation} passent-ils avant
les fonds des survivants et non après~?
\item Un membre soutient que les contributions supplémentaires ne devraient pas être plafonnées «~pour que la
chambre de \omterm{def:m1:clearing-and-settlement:clearing}{compensation} ne puisse jamais faire défaut~». Donner l'argument contraire.
\item Le même lundi, chacun des autres membres a reçu un important
appel de marge intrajournalier. Pourquoi cela constitue-t-il un risque pour le système, et comment
appelle-t-on ce phénomène~?
\item Énoncer le \emph{résultat nommé}~: la perte atteignant le \omterm{def:m1:clearing-and-settlement:waterfall}{fonds de
défaillance} mutualisé, en millions de dollars.
\item En une phrase~: qui garantit, en définitive, une transaction compensée~?
\end{enumerate}
\end{problem}

\section{Questions d'entretien}

\begin{interviewq}[$\star$ \iqroles{trader, développeur, banque}]\label{iq:m1:clearing-and-settlement:1}
Qu'est-ce que la \omterm{def:m1:clearing-and-settlement:ccp}{novation}, et quel risque supprime-t-elle pour une société de trading~?
\end{interviewq}

\begin{interviewq}[$\star$ \iqroles{trader, chercheur, banque}]\label{iq:m1:clearing-and-settlement:2}
Quelle est la différence entre la \omterm{def:m1:clearing-and-settlement:margin}{marge initiale} et la \omterm{def:m1:clearing-and-settlement:margin}{marge de variation}~? Laquelle
des deux vous est restituée~?
\end{interviewq}

\begin{interviewq}[$\star\star$ \iqroles{trader, chercheur}]\label{iq:m1:clearing-and-settlement:3}
En janvier 2021, plusieurs courtiers de détail ont empêché leurs clients
d'acheter certaines actions. Donnez l'explication mécanique.
\end{interviewq}

\begin{interviewq}[$\star\star$ \iqroles{banque, chercheur}]\label{iq:m1:clearing-and-settlement:4}
Décrivez la \omterm{def:m1:clearing-and-settlement:waterfall}{cascade de défaillance} d'une chambre de \omterm{def:m1:clearing-and-settlement:clearing}{compensation}. Pourquoi cet
ordre est-il tel qu'il est~?
\end{interviewq}

\begin{interviewq}[$\star\star$ \iqroles{développeur}]\label{iq:m1:clearing-and-settlement:5}
Vous rédigez le traitement de netting de fin de journée. Quels invariants vérifiez-vous
avant de publier les instructions de \omterm{def:m1:clearing-and-settlement:clearing}{règlement-livraison}, et que faites-vous si l'un d'eux
échoue à 18:55, alors que l'échéance est fixée à 19:00~?
\end{interviewq}

\begin{interviewq}[$\star\star\star$ \iqroles{chercheur, banque}]\label{iq:m1:clearing-and-settlement:6}
La \omterm{def:m1:clearing-and-settlement:clearing}{compensation} centrale a été rendue obligatoire pour les swaps standard après 2008.
Réduit-elle le risque systémique~? Argumenter les deux positions.
\end{interviewq}
